% !PRINTED 13
\weberpage{13}

%% [running head, not body text -- commented out 2026-08-07] Page furniture (running
%% head, folio, signature) is excluded from this edition by standing policy.
%% \gesperrt{\S.\,4.} \hfill \gesperrt{Group of the subsidiary groups.} \hfill 13

\gesperrt{2. By the two equations}

\weq{2}{AA = A,}
\weq{3}{AB = BA,}

\gesperrt{in which $A$ may be a definite part, $B$ any arbitrary part of $P$, it is asserted that $A$ is a normal divisor of $P$.}

For by (2), $A$ is a divisor of $P$, and by (3), for every element $b$ of $P$
\[
b^{-1}Ab = A,
\]
i.\,e. $A$ is a normal divisor.

\gesperrt{3. If $A$ and $B$ are two divisors of $P$ (groups), then either of the two equations}

\weq{4}{AB = A, \quad BA = A,}

\gesperrt{is the necessary and sufficient condition for $B$ to be a divisor of $A$.}

For if $B$ is a group and a divisor of the group $A$, and if $a$ is contained in $A$, and $b$ in $B$ and hence also in $A$, then $ab$ is also contained in $A$. Thus $A$ contains every element of $AB$.

But since, secondly, $B$ as a group also contains the element 1, $AB$ also contains every element of $A$, and hence $AB$ is identical with $A$.

In the same way one sees that $BA$ is identical with $A$. Conversely, it follows from either of the equations (4), since $A$ as a group contains the unit element, that every element of $B$ is contained in $A$, and hence that $B$ is a divisor of $A$.

\gesperrt{4. Under the composition of the parts, the system of the subsidiary groups belonging to a normal divisor of $P$ itself forms a group, in which the normal divisor $N$ forms the unit, and}
\[
Nb, \quad Nb^{-1}
\]
\gesperrt{are opposite elements.}

For if $N$ is a normal divisor of $P$ and
\weq{5}{N, \; N_1 = Nb_1, \; N_2 = Nb_2, \dots}
is the system of the subsidiary groups, then
\[
N_h N_k = N b_h N b_k.
\]
